\documentclass[11pt]{article}

\usepackage[margin=1in]{geometry}
\usepackage{amsmath,amssymb,amsthm,mathtools}
\usepackage{array}
\usepackage{booktabs}
\usepackage{enumitem}
\usepackage{graphicx}
\usepackage{microtype}
\usepackage{url}
\usepackage[numbers]{natbib}
\usepackage[colorlinks=true,citecolor=blue,linkcolor=blue,urlcolor=blue]{hyperref}

\newtheorem{proposition}{Proposition}
\newtheorem{assumption}{Assumption}
\newtheorem{definition}{Definition}
\newcolumntype{L}[1]{>{\raggedright\arraybackslash}p{#1}}
\newcommand{\ind}{\mathbb{1}}
\newcommand{\calD}{\mathcal{D}}
\newcommand{\norm}[1]{\left\lVert #1 \right\rVert}
\newcommand{\TODOFIG}[2]{%
\begin{center}
\fbox{\begin{minipage}{0.88\linewidth}\vspace{0.6em}
\textbf{Figure placeholder: #1}\\[0.35em]
#2
\vspace{0.6em}\end{minipage}}
\end{center}}

\setlength{\textfloatsep}{8pt plus 2pt minus 2pt}
\setlength{\floatsep}{6pt plus 2pt minus 2pt}
\setlength{\intextsep}{8pt plus 2pt minus 2pt}
\setlength{\abovecaptionskip}{3pt plus 1pt minus 1pt}
\setlength{\belowcaptionskip}{2pt plus 1pt minus 1pt}

\title{Budgeted Max-Min Transitive Reachability for Long-Horizon Offline Goal-Conditioned Control}
\author{Anonymous Authors}
\date{}

\begin{document}

\maketitle

\begin{abstract}
Offline goal-conditioned reinforcement learning asks an agent to compose short
behavioral fragments in a static dataset into policies that reach distant goals.
This is a fundamental long-horizon credit-assignment problem: online correction
is unavailable, value errors may be amplified by bootstrapping, and unsupported
intermediate goals can derail execution. Transitive RL (TRL) recently proposed a
promising divide-and-conquer value update through intermediate goals, reducing
recursive dependency depth. However, the discounted product backup used by TRL is
addition in temporal-distance space, so approximation errors from the two
branches can still add even when the recursion tree is shallow. We propose
Budgeted Max-Min Transitive RL (BMM-TRL), which learns finite-budget
reachability predicates and composes them with a max-min backup. For deterministic
or support-graph reachability, max-min composition is non-expansive in sup norm,
yielding an ideal max-branch-error recurrence and logarithmic accumulated score
error under balanced recursion. Empirically, BMM improves heldout long-budget
Q/reachability classification on grid-cell, environment-step, and Scene-Play
support-graph diagnostics, with product-style transitive targets serving as a
strong but slightly weaker control. When paired with fixed local goal-conditioned
controllers, BMM reachability scores also provide an effective high-level
subgoal interface on long-horizon OGBench-style tasks, including several
reported TRL benchmark rows. The resulting picture is deliberately modular:
BMM-TRL addresses long-horizon value propagation; fixed-controller planning
experiments show one practical policy interface; end-to-end actor extraction from
BMM values remains an important open problem.
\end{abstract}

\section{Introduction}

Goal-conditioned reinforcement learning (GCRL) aims to learn policies
$\pi(a\mid s,g)$ that can reach many goals from many states
\citep{kaelbling1993learning,schaul2015uvfa,andrychowicz2017her}. In the offline
setting, this promise becomes especially attractive: large reward-free datasets
could be converted into broadly reusable behaviors without online interaction
\citep{fu2020d4rl,park2024ogbench}. Yet the same setting exposes a central
failure mode. Long-horizon success requires composing local transitions into
far-away reachability estimates, but offline agents cannot collect data to repair
value mistakes, and learned policies must avoid actions and subgoals outside the
data support \citep{kumar2020cql,kostrikov2022iql}.

The question we study is therefore simple:
\begin{quote}
\emph{What algebra should a learned goal-conditioned value use when it composes
short-horizon reachability into long-horizon reachability?}
\end{quote}
The answer matters because the algebra controls how approximation errors
propagate. Standard temporal-difference backups propagate information one step at
a time \citep{sutton1988td,watkins1992qlearning}. Transitive RL (TRL) introduced
an appealing alternative: use an intermediate witness $w$ and compose values for
$s\to w$ and $w\to g$ \citep{park2025trl}. In idealized settings this can reduce
the number of recursive compositions needed to reason about a length-$T$ path
from $O(T)$ to $O(\log T)$. This is an important idea, and we keep it.

Our starting observation is that recursion depth and numerical error algebra are
not the same. In deterministic goal reaching, discounted values satisfy
\begin{equation}
    V^*(s,g)=\gamma^{d^*(s,g)},
\end{equation}
where $d^*(s,g)$ is shortest-path distance. TRL's product backup
\begin{equation}
    V(s,g) \leftarrow \max_w V(s,w)V(w,g)
\end{equation}
is equivalent, after taking logarithms, to additive distance composition,
\begin{equation}
    d(s,g) \leftarrow \min_w d(s,w)+d(w,g).
\end{equation}
Thus errors in two balanced branches can still add in distance space. A shallow
computation graph alone does not guarantee a shallow error recurrence.

BMM-TRL changes the object being learned. Instead of estimating an exponentiated
exact distance, it estimates finite-budget reachability,
\begin{equation}
    R_H(s,g)=\ind\{d^*(s,g)\le H\},
\end{equation}
and composes it through
\begin{equation}
    R_H(s,g)=\max_w \min\big(R_h(s,w),R_{H-h}(w,g)\big).
\end{equation}
The max and min operators are non-expansive under sup-norm perturbations. Hence,
if the two child predictors have errors $E_h$ and $E_{H-h}$ and the fitted parent
incurs residual $\epsilon_H$, the ideal recurrence is
\begin{equation}
    E_H \le \epsilon_H + \max(E_h,E_{H-h}),
\end{equation}
not $\epsilon_H+E_h+E_{H-h}$. Under balanced dyadic budgets this yields
$O(\epsilon\log H)$ accumulated score error for uniform residual $\epsilon$.

This paper makes three contributions.
First, we formulate BMM-TRL, a budgeted reachability version of transitive value
learning whose backup is max-min rather than product. Second, we give a simple
non-expansive error-propagation result for deterministic and support-graph
reachability, directly addressing the gap between shallow recursion and additive
branch error. Third, we evaluate the idea with a focused experimental design: a
budget-holdout diagnostic isolates long-budget value bootstrapping; product
controls test whether the max-min algebra matters beyond transitive transfer
itself; and fixed-controller planning experiments measure whether the learned
reachability scores can guide long-horizon behavior while keeping policy
extraction separate.

\begin{figure}[t]
\centering
\TODOFIG{Core motivation and method schematic}{Draw three panels. Left: TD propagates value through a length-$T$ chain. Middle: TRL uses a balanced witness tree but product composition becomes addition in distance space, so branch errors add. Right: BMM learns budgeted reachability heads and composes them with max-min, so branch errors propagate by maximum. Include the recurrences $E_T=E_{T-1}+\epsilon$, $E_T=2E_{T/2}+\epsilon$, and $E_T=E_{T/2}+\epsilon$.}
\caption{BMM-TRL keeps the transitive witness idea but changes the value algebra.}
\label{fig:motivation}
\end{figure}

\section{Problem Setting and Background}

\paragraph{Offline goal-conditioned control.}
We assume a fixed dataset
$\calD=\{(s_t,a_t,s_{t+1})\}$ collected without access to the downstream test
goals. A goal-conditioned agent receives a current state $s$ and goal $g$ and
must output actions that reach $g$. We focus on the long-horizon regime where
successful behavior requires stitching many local transitions. In this regime,
value learning and policy extraction are distinct bottlenecks: a value may fail
to infer distant reachability, and a policy may fail to convert even a good value
into supported actions.

\paragraph{Transitive value learning.}
TRL composes goal-conditioned values through intermediate states
\citep{park2025trl}. Its practical target has the form
\begin{equation}
    Q(s,a,g) \leftarrow \max_{w,a'} Q(s,a,w)Q(w,a',g),
\end{equation}
with in-trajectory witnesses used to reduce overestimation. This is a powerful
way to exploit the triangle-inequality structure of goal reaching. Our critique
is narrower: product composition is the correct algebra for exponentiated
additive distance, but it does not by itself give a max-error recurrence in
log-value or distance space.

\paragraph{Support reachability.}
In offline control, the relevant question is often not only whether a geometric
path exists, but whether it exists through the dataset support. We therefore use
three reachability notions in experiments: grid/geodesic reachability for clean
maze diagnostics, calibrated environment-step reachability, and support-graph
reachability in learned or oracle representations. The theory below applies
exactly to deterministic shortest-path graphs and to any fixed support graph;
stochastic extensions require additional reliability semantics and are discussed
in Section~\ref{sec:limitations}.

\section{Budgeted Max-Min Transitive RL}

\subsection{Budgeted reachability values}

For a deterministic graph or deterministic MDP, let $d^*(s,g)$ denote the
shortest supported path length. Define the state reachability predicate
\begin{equation}
    V_H^*(s,g)=\ind\{d^*(s,g)\le H\}.
\end{equation}
For an action-conditioned critic with deterministic transition $s'=f(s,a)$, we
use
\begin{equation}
    Q_H^*(s,a,g)=\ind\{d^*(s',g)\le H-1\}.
\end{equation}
A neural critic outputs a logit $f_\theta(s,a,g,H)$ and reachability score
\begin{equation}
    R_\theta(s,a,g,H)=\sigma(f_\theta(s,a,g,H)).
\end{equation}
The budget $H$ is appended to the goal input using a normalized log-budget
feature, optionally with a one-hot budget code for multi-budget diagnostics.

\subsection{Exact max-min identity}

\begin{proposition}[Budgeted transitivity]
For any split $h\in\{0,\ldots,H\}$, allowing endpoint witnesses,
\begin{equation}
    V_H^*(s,g)=\max_w \min\left(V_h^*(s,w),V_{H-h}^*(w,g)\right).
\end{equation}
The action-conditioned analogue is
\begin{equation}
    Q_H^*(s,a,g)=\max_w \min\left(Q_h^*(s,a,w),V_{H-h}^*(w,g)\right).
\end{equation}
\end{proposition}

\begin{proof}
If $V_H^*(s,g)=1$, then a path of length at most $H$ exists. Choose the witness
$w$ reached after at most $h$ steps along that path; both branch predicates are
one. Conversely, if some witness makes both branch predicates one, concatenating
the two supported paths gives a supported path from $s$ to $g$ of length at most
$H$. The action-conditioned identity applies the same argument after the first
action transition.
\end{proof}

\subsection{Non-expansive error propagation}

\begin{proposition}[Max-branch error propagation]
Suppose estimates $\widehat V_h$ and $\widehat V_{H-h}$ satisfy
\begin{equation}
    \norm{\widehat V_h - V_h^*}_\infty \le E_h,
    \qquad
    \norm{\widehat V_{H-h} - V_{H-h}^*}_\infty \le E_{H-h}.
\end{equation}
Define
\begin{equation}
    \widehat T_H(s,g)=\max_w \min\left(\widehat V_h(s,w),\widehat V_{H-h}(w,g)\right).
\end{equation}
Then
\begin{equation}
    \norm{\widehat T_H-T_H^*}_\infty\le \max(E_h,E_{H-h}).
\end{equation}
\end{proposition}

\begin{proof}
For scalars, $|\min(a,b)-\min(a',b')|\le
\max(|a-a'|,|b-b'|)$. The maximum over witnesses is also non-expansive:
$|\max_w z_w-\max_w z'_w|\le \sup_w |z_w-z'_w|$. Combining these two inequalities
gives the claim.
\end{proof}

If fitting the parent predictor introduces residual $\epsilon_H$, then
\begin{equation}
    E_H \le \epsilon_H + \max(E_h,E_{H-h}).
\end{equation}
For dyadic budgets with balanced splits and uniform residual $\epsilon$,
$E_H\le \epsilon\log_2 H$. By contrast, additive distance composition obeys
\begin{equation}
    E_H \le \epsilon_H+E_h+E_{H-h},
\end{equation}
which is linear in horizon under equal splits and uniform residuals. This is the
central algebraic difference between BMM and product-style transitive learning.

\subsection{Training objective}

The policy-relevant diagnostic uses an action-conditioned Q critic and a frozen
state-value teacher for the second branch. For a parent budget $H$ and split
$(h,H-h)$, the BMM target is
\begin{equation}
    y_{\mathrm{BMM}}(s,a,g,H)=
    \max_{w\in\mathcal{W}} \min\left(Q_h(s,a,w),V_{H-h}(w,g)\right),
    \label{eq:bmm-qv}
\end{equation}
where $\mathcal{W}$ is a finite set of supported witnesses. Because finite
witness sampling gives a lower bound on the true maximum, we use conservative
binary cross-entropy targets and mask invalid witness samples. The matched
product control replaces the min in Eq.~\ref{eq:bmm-qv} by multiplication:
\begin{equation}
    y_{\mathrm{prod}}(s,a,g,H)=
    \max_{w\in\mathcal{W}} Q_h(s,a,w)V_{H-h}(w,g).
\end{equation}
We also compare to a no-parent supervised baseline and a V-next distillation
control that does not perform transitive Q/V composition.

\paragraph{Policy interface.}
BMM does not require a new low-level actor API. In our control experiments, we
fix a local goal-conditioned controller and use BMM scores only for high-level
subgoal selection. This design isolates the question ``are BMM values useful for
choosing supported intermediate goals?'' from the separate question ``can a flat
actor be extracted directly from BMM Q-values?'' The latter remains open.

\begin{figure}[t]
\centering
\TODOFIG{Budget-holdout diagnostic}{Show short-budget supervision heads $H/2$ and heldout parent head $H$. Direct parent labels are hidden from the transitive variant; the parent is learned through sampled witnesses and the $Q_h,V_{H-h}$ target. Add a small inset comparing no-transitive, product, max-min, and V-next.}
\caption{Budget holdout isolates long-budget value bootstrapping from policy execution.}
\label{fig:budget-holdout-design}
\end{figure}

\section{Experimental Design}

The experiments are organized around three questions.
\begin{enumerate}[leftmargin=1.4em]
    \item \textbf{Does the backup algebra change error scaling?} We run a
    tabular recurrence sanity check with controlled residuals.
    \item \textbf{Can short-budget knowledge bootstrap a heldout long-budget
    reachability predictor?} We train short-budget Q/V predictors and withhold
    direct parent labels, comparing BMM, product, no-transitive, and V-next
    controls.
    \item \textbf{Are BMM reachability scores useful for long-horizon behavior
    when policy extraction is controlled?} We use fixed local controllers and
    compare high-level subgoal interfaces on OGBench-style tasks.
\end{enumerate}

\paragraph{Value diagnostics.}
Budget-holdout is the main value-learning test. The model receives short-budget
supervision and must learn a parent budget through transitive Q/V targets. We
report AUC, positive-negative score gap, and binary cross-entropy (BCE); higher
AUC/gap and lower BCE are better. The no-transitive baseline tests whether the
parent can be recovered without a transitive target, the V-next baseline tests a
non-compositional teacher, and the product baseline tests whether transitive
transfer alone explains the gains.

\paragraph{Control protocol and attribution.}
Our policy-facing experiments are fixed-controller hierarchical evaluations.
The low-level controller is held fixed within each comparison, while the
high-level subgoal rule varies among geometric, support-only, product/BMM value,
and calibrated BMM/support selectors when available. This is a system-level
comparison to reported benchmark rows, not an apples-to-apples flat-policy
replacement for TRL's RPG or rejection-sampling policy extraction. We make this
scope explicit because it is also a scientific point: long-horizon offline GCRL
has a value-propagation bottleneck and a policy-interface bottleneck.

\section{Results}

\subsection{Max-min gives the intended error recurrence}

The tabular recurrence check injects a uniform per-level residual
$\epsilon=0.02$ and compares balanced max-min composition to additive/product
composition. At horizon $H=1024$, the BMM recurrence gives error $0.22$, whereas
additive/product composition gives $20.48$. This experiment is intentionally
simple: its purpose is to verify the algebraic prediction in a setting where the
only difference is the composition rule.

\begin{figure}[t]
\centering
\TODOFIG{Error-scaling curve}{Plot horizon on the x-axis, sup error on the y-axis. Use a log-scale x-axis. Show BMM balanced error growing as $O(\log H)$ and additive/product error growing as $O(H)$. Annotate $H=1024$: 0.22 vs 20.48.}
\caption{Controlled recurrence sanity check. Max-min composition propagates the larger child error; product/distance composition adds branch errors.}
\label{fig:error-scaling}
\end{figure}

\subsection{Budget-holdout value learning}

Table~\ref{tab:budget-holdout} shows the central value-learning result. BMM
improves heldout parent-budget classification over no-transitive baselines, while
V-next is near zero. Product targets also transfer, but BMM has a modest and
consistent edge in these matched diagnostics.

\begin{table}[t]
\centering
\caption{Heldout parent-budget Q classification. Deltas are relative to the named baseline; higher $\Delta$AUC/$\Delta$gap and lower $\Delta$BCE are better.}
\label{tab:budget-holdout}
\begin{tabular}{llrrr}
\toprule
Setting & Comparison & $\Delta$AUC & $\Delta$gap & $\Delta$BCE \\
\midrule
grid-cell H8 & BMM $-$ no trans & +0.0150 & +0.0751 & -0.3609 \\
grid-cell H8 & product $-$ no trans & +0.0139 & +0.0702 & -0.3344 \\
grid-cell H8 & BMM $-$ product & +0.0011 & +0.0050 & -0.0265 \\
grid-cell H8 & V-next $-$ no trans & +0.0002 & +0.0017 & -0.0063 \\
\midrule
env-step H160 & BMM $-$ no trans & +0.0243 & +0.0647 & -0.5815 \\
env-step H160 & product $-$ no trans & +0.0215 & +0.0539 & -0.5037 \\
env-step H160 & BMM $-$ product & +0.0028 & +0.0108 & -0.0778 \\
env-step H160 & V-next $-$ no trans & +0.0035 & +0.0030 & -0.0570 \\
\bottomrule
\end{tabular}
\end{table}

\paragraph{Non-maze support-graph transfer.}
To test whether the value-learning effect depends on maze coordinates, we build
support graphs over Scene-Play oracle representations. Table~\ref{tab:scene-qv}
shows that BMM and product Q/V targets substantially improve H128 reachability
classification over no-transitive and V-next controls. The train-only graph uses
only train-support bins; thus the result tests transfer through the offline
support representation rather than direct access to validation transitions.

\begin{table}[t]
\centering
\caption{Scene-Play H128 support-graph Q/V holdout. These are value diagnostics, not policy-control results.}
\label{tab:scene-qv}
\begin{tabular}{llrr}
\toprule
Setting & Variant & AUC & gap \\
\midrule
train+val graph & no trans & 0.6664 & 0.0065 \\
train+val graph & BMM Q/V & 0.8177 & 0.2355 \\
train+val graph & product Q/V & 0.8129 & 0.1969 \\
train+val graph & V-next & 0.6694 & 0.0067 \\
\midrule
train-only graph & no trans & 0.6171 & 0.0128 \\
train-only graph & BMM Q/V & 0.7618 & 0.1896 \\
train-only graph & product Q/V & 0.7568 & 0.1563 \\
train-only graph & V-next & 0.6181 & 0.0130 \\
\bottomrule
\end{tabular}
\end{table}

\subsection{Fixed-controller long-horizon control}

BMM values are useful when used as a high-level subgoal scorer for a fixed local
controller. Table~\ref{tab:fixed-actor} reports rows where the same general
low-level controller family is held fixed and BMM changes the high-level
interface. These rows are the cleanest policy-facing evidence because they test
whether BMM subgoal selection improves long-horizon execution beyond direct or
support-only controls.

\begin{table}[t]
\centering
\small
\caption{Fixed-controller graph-subgoal results. ``Reported TRL'' is the success rate listed in the TRL paper table; our protocol adds a high-level BMM/support subgoal interface while keeping a fixed local controller.}
\label{tab:fixed-actor}
\begin{tabular}{L{0.32\linewidth}rrL{0.34\linewidth}}
\toprule
Environment & Reported TRL & Ours & Protocol / attribution \\
\midrule
humanoidmaze-medium-navigate-oraclerep & 57.0\% & 94.7\% (71/75) & BMM support-path graph subgoals; official 2000-step horizon \\
humanoidmaze-large-navigate-oraclerep & 8.0\% & 89.3\% (67/75) & BMM support-path graph subgoals; official 2000-step horizon \\
humanoidmaze-giant-navigate-oraclerep & 79.0\% & 82.7\% (62/75) & calibrated BMM/support route selector; pure selectors remain below the oracle union \\
antsoccer-arena-navigate-oraclerep & 73.0\% & 90.7\% (68/75) & same 1M TRL/RPG actor; BMM subgoal selection; direct actor 53/75, support control 59/75 \\
scene-play-oraclerep & 77.0\% & 86.7\% (65/75) & train-only support graph plus local GCFBC; matched support-only 63/75 \\
\bottomrule
\end{tabular}
\end{table}

\begin{figure}[t]
\centering
\TODOFIG{Hierarchical policy interface}{Draw a start and final goal connected by a support graph. Show support-only choosing a shortest support path, product/BMM scoring candidate witnesses, and a fixed low-level controller executing each selected subgoal. Include an annotation: ``same local controller, different high-level selector''.}
\caption{Policy-facing experiments evaluate BMM as a value-guided subgoal interface, not as a new flat actor extractor.}
\label{fig:policy-interface}
\end{figure}

\paragraph{Structured-interface case studies.}
Several hard manipulation and puzzle rows are best understood as policy-interface
case studies rather than pure BMM value-ranking results. Table~\ref{tab:structured}
shows that exact or structured decompositions can dramatically improve success
with learned local controllers. These results are useful because they show that
some reported long-horizon failures are not due to irrecoverable dataset
coverage; however, they should not be attributed solely to the BMM backup.

\begin{table}[t]
\centering
\small
\caption{Structured policy-interface case studies. These rows compare against reported TRL success rates but use task-structured or local-controller interfaces.}
\label{tab:structured}
\begin{tabular}{L{0.34\linewidth}rrL{0.30\linewidth}}
\toprule
Environment & Reported TRL & Ours & Interface \\
\midrule
puzzle-4x4-play-oraclerep & 34.0\% & 97.3\% (73/75) & GF(2) Lights-Out solve + local controller \\
puzzle-4x5-play-oraclerep & 97.0\% & 100.0\% (75/75) & exact discrete planner + local controller \\
puzzle-4x6-play-oraclerep & 51.0\% & 92.0\% (69/75) & exact discrete planner + local controller \\
cube-single-play-oraclerep & 95.0\% & 98.7\% (74/75) & direct local GCFBC controller \\
cube-double-play-oraclerep & 30.0\% & 73.3\% (55/75) & dynamic sequential block subgoals + local controller \\
\bottomrule
\end{tabular}
\end{table}

\subsection{What the experiments establish}

The strongest empirical claim is about value learning: BMM improves heldout
long-budget Q/reachability prediction without direct long-budget labels. The
strongest policy claim is conditional: BMM scores help a fixed local controller
choose useful subgoals in long-horizon tasks. The experiments do not establish
that BMM is a drop-in replacement for TRL's flat RPG or rejection-sampling policy
extraction. This distinction is central to the paper's interpretation, not a
minor caveat.

\section{Related Work}

BMM-TRL builds on several lines of work. Classical dynamic programming and
shortest-path algorithms establish the value-composition view of sequential
decision making \citep{bellman1957dynamic,bertsekas2017dp}. Temporal-difference
and Q-learning methods propagate value locally \citep{sutton1988td,watkins1992qlearning},
while modern offline RL studies how to learn from fixed datasets without
extrapolating to unsupported actions \citep{fujimoto2019bcq,kumar2019bear,kumar2020cql,fujimoto2021td3bc,kostrikov2022iql}.
Goal-conditioned RL has a long history in universal value functions and
hindsight relabeling \citep{kaelbling1993learning,schaul2015uvfa,andrychowicz2017her},
with recent offline GCRL algorithms using contrastive objectives, hierarchical
latent actions, and benchmark-driven evaluation
\citep{eysenbach2022contrastive,park2023hiql,park2024ogbench,zhu2023provably}.
TRL is the closest predecessor: it uses witness-based transitive value learning
for long-horizon offline GCRL \citep{park2025trl}. BMM keeps the witness idea but
changes the algebra from product to max-min finite-budget reachability.

Our policy interface is related to hierarchical and graph-based planning with
learned local controllers. Prior work plans over replay-buffer or topological
memory graphs \citep{savinov2018sptm,eysenbach2019sorb} and learns hierarchical
controllers through subgoals \citep{nachum2018hiro,levy2019learning,nair2018rig}.
Recent work also studies test-time graph search and representation-aligned
planning for offline GCRL \citep{opryshko2025ttgs,kang2026lavl}. BMM differs in
that its main contribution is not the existence of a graph planner, but a
non-expansive transitive value target that can score finite-budget reachability
on such support structures.

\section{Limitations and Future Work}
\label{sec:limitations}

\paragraph{Stochastic dynamics.}
The max-min identity is exact for deterministic shortest paths and fixed support
graphs. In stochastic domains, raw success probability does not generally obey a
hard max-min identity. A natural extension is to learn reliability-threshold
reachability, e.g., whether the probability of reaching $g$ within $H$ steps
exceeds a confidence level. This may preserve a conservative bottleneck
interpretation, but it requires separate analysis.

\paragraph{Support graph construction.}
Current support graphs depend on representation choice, binning, edge rules, and
horizon scaling. Better automatic graph construction, learned representations,
and uncertainty-aware support estimates are important for broader deployment.

\paragraph{Policy extraction.}
BMM's direct actor interface is not solved. Flat action ranking, direct
RPG/FRS-style extraction, and joint action-subgoal extraction were weaker than
fixed-controller planning in our current experiments. Future work should combine
BMM values with standardized behavior-constrained actor optimization and
rejection-sampling policies under the same protocol used by TRL.

\paragraph{Benchmark comparability.}
The benchmark tables compare to reported TRL success rates, but many of our
policy rows add a high-level subgoal interface. This is a legitimate system-level
comparison, but not an apples-to-apples comparison of flat policy extraction. We
therefore report protocol qualifiers and matched support-only/direct-controller
controls where available.

\section{Conclusion}

BMM-TRL reframes transitive value learning around finite-budget reachability.
This changes the composition rule from product/additive distance composition to
max-min reachability composition, giving a non-expansive error-propagation
operator in deterministic and support-graph settings. Budget-holdout experiments
show that this target improves heldout long-budget Q/reachability prediction,
and fixed-controller evaluations show that the resulting values can guide
long-horizon subgoal planning. The broader message is that long-horizon offline
GCRL requires both reliable value composition and effective policy interfaces;
BMM addresses the former and provides a useful tool for the latter.

\appendix

\section{Additional Proof Details}

\subsection{Balanced dyadic recursion}

Let $H=2^L$ and suppose each parent fit introduces residual at most $\epsilon$.
Repeatedly applying Proposition~2 with balanced splits gives
\begin{align}
E_{2^L} &\le \epsilon + E_{2^{L-1}} \\
&\le 2\epsilon + E_{2^{L-2}} \\
&\le \cdots \le L\epsilon + E_1.
\end{align}
If the base budget is fit exactly, $E_1=0$ and $E_H\le \epsilon\log_2 H$. For
additive distance composition,
$E_{2^L}\le \epsilon+2E_{2^{L-1}}$, yielding
$E_{2^L}\le (2^L-1)\epsilon$ when $E_1=0$.

\subsection{Why product and max-min can both transfer}

Product Q/V targets can improve heldout reachability because they also provide a
transitive signal through witnesses. The difference is in worst-case score-error
algebra. Product is the right composition for exponentiated additive distance;
max-min is the right composition for binary finite-budget reachability. Thus our
experiments treat product as a strong control rather than a strawman baseline.

\section{Implementation Details}

\paragraph{Budget features.}
Budgets are represented by $\log_2(H)/\log_2(H_{\max})$, optionally concatenated
with a one-hot code over the active budget set. One-hot budget features were used
in the harder multi-budget holdout diagnostics to reduce interpolation
ambiguity.

\paragraph{Witness sampling.}
For budget-holdout Q/V experiments, witnesses are sampled from supported states
or same-trajectory midpoints depending on the target type. Invalid or degenerate
splits are masked. The target networks are frozen for the branch scores when
forming parent targets.

\paragraph{Support-only and calibrated selectors.}
A support-only selector follows a path in the offline support graph without using
BMM as the primary score. A calibrated selector chooses between BMM-guided and
support-only routes using a small fixed gate over start-goal features such as
support-graph distance or displacement. These selectors are reported explicitly
because they are stronger policy interfaces than flat actor extraction.

\section{Additional Experimental Tables}

\begin{table}[h]
\centering
\caption{Claim scope and evidence map.}
\label{tab:claim-scope}
\begin{tabular}{L{0.24\linewidth}L{0.42\linewidth}L{0.24\linewidth}}
\toprule
Claim type & Evidence & Paper use \\
\midrule
Value-error propagation & non-expansive proof and recurrence sanity check & main theoretical claim \\
Heldout value learning & grid-cell, env-step, Scene-Play support-graph holdouts & main empirical claim \\
BMM vs product & matched product controls & product is strong; BMM edge is modest \\
Long-horizon behavior & fixed-controller subgoal selection & policy-facing evidence \\
Flat policy extraction & current direct extraction weak & limitation / future work \\
\bottomrule
\end{tabular}
\end{table}

\begin{table}[h]
\centering
\small
\caption{Representative fixed-controller controls from development runs. These rows are useful for attribution but should be expanded before camera-ready submission.}
\label{tab:controls}
\begin{tabular}{L{0.32\linewidth}L{0.28\linewidth}L{0.30\linewidth}}
\toprule
Setting & Baseline/control & BMM-side result \\
\midrule
medium PointMaze navigate & geometric 40.0\% & support-gated BMM 96.0\% \\
medium PointMaze stitch & geometric 20.0\% & support-gated BMM 100.0\% \\
large PointMaze stitch & geometric 0.0\%; support-only 20.0\% & value-frontier BMM 100.0\% \\
antmaze-medium & support-only 93.3\% & BMM support-path 100.0\% \\
antmaze-large-oraclerep & support-only 80.0\% & support-path plus BMM repair 86.7\% \\
\bottomrule
\end{tabular}
\end{table}

\section{Reproducibility Notes}

All headline tables are regenerated from local JSON artifacts by the claim
validator:
\begin{verbatim}
conda run --no-capture-output -n bmm-trl python \
  scripts/validate_bmm_paper_claims.py
\end{verbatim}
The exact rollout commands for representative Scene-Play, Puzzle, Cube, and
AntSoccer evaluations are listed in \path{BMM_TRL_REPRO_COMMANDS.md}. Before
submission, we recommend freezing a single evaluation protocol per environment,
committing the generated JSON summaries, and adding a one-command artifact check
for every number in the main tables.

\section{Figure Checklist for Camera Ready}

\begin{enumerate}[leftmargin=1.4em]
    \item \textbf{Motivation schematic}: TD chain vs TRL product tree vs BMM
    max-min tree.
    \item \textbf{Error-scaling curve}: controlled recurrence with $H$ on a log
    x-axis.
    \item \textbf{Budget-holdout diagram}: short-budget supervision, hidden
    parent budget, and transitive target construction.
    \item \textbf{Policy-interface diagram}: support graph, candidate witnesses,
    BMM scoring, and fixed local controller execution.
\end{enumerate}

\bibliographystyle{plainnat}
\bibliography{references}

\end{document}
